% ARCHIVED DRAFT: independently audited by root and incorporated on 2026-11-19 16:21 UTC.
% The assembled paper is authoritative; this file retains drafting provenance.
\begin{saresult}{B.068}{Proved under stated hypotheses}{Regular Vandermonde ideals and arbitrary coefficient extensions}
Let $R$ be a nonzero commutative ring, let $Q$ be its total quotient ring,
and let $D_d(R)$ be the Vandermonde ideals of B.064. Suppose that, for every
$d\ge1$, the ideal $D_d(R)$ contains a nonzerodivisor of $R$, and
\[
                         (R:_Q D_d(R))=R.                 \tag{B68.1}
\]
Then for every unital inclusion $R\subset S$ the following are equivalent:
\[
\begin{gathered}
 \text{the polynomial co-operations exist uniquely};\\
 \theta_2\text{ is injective};\\
 \Omega(R,S)\longrightarrow\operatorname{App}(R,R)
       \text{ is injective};\\
 \operatorname{Ann}_S(D_d(R))=0\quad\text{for every }d;\\
 \Omega(R,S)=R[X].                                         \tag{B68.2}
\end{gathered}
\]
Every positive case has all multivariable polynomial tensor isomorphisms.
The hypotheses hold for every infinite domain with
$\operatorname{Int}(R,\operatorname{Frac}R)=R[X]$.
They also hold for $R=A[Z,W]$ over every nonzero commutative ring $A$,
without finiteness or reducedness assumptions. More generally they hold
for $R=A[Z]$ whenever each $D_d(A)$ contains a nonzerodivisor; in particular,
this includes every infinite domain $A$.
\end{saresult}
\begin{proof}
Write $U$ for the multiplicative set of nonzerodivisors in $R$.
The inclusion $R\subset S$ induces an inclusion $Q\subset U^{-1}S$:
if an element of $R$ becomes zero, a nonzerodivisor of $R$ annihilates it,
which is impossible unless it was zero.
For $F=\sum_{i=0}^d a_iX^i\in\Omega(R,S)$, Cramer's rule gives
$D_d(R)a_i\subset R$ for each coefficient. Choose a nonzerodivisor
$v\in D_d(R)$; it may be a finite linear combination of Vandermonde
products, since the Cramer conclusion holds on their whole ideal. The image of $a_i$ in $U^{-1}S$ equals
$(va_i)/v\in Q$. The rationalized polynomial is therefore in
$\operatorname{Int}(R,Q)$. By B.064 and (B68.1),
$\operatorname{Int}(R,Q)=R[X]$.
Rationalization consequently gives an $R$-linear retraction
\[
 \rho:\Omega(R,S)\longrightarrow R[X],\qquad
 \Omega(R,S)=R[X]\oplus\mathcal N,                         \tag{B68.3}
\]
where $\mathcal N$ is exactly the module of null polynomials on $R$.
Indeed an element in the kernel has each coefficient killed by some
member of $U$, hence all its coefficients by a common such member.
Its values belong to $R$, so regularity makes all values zero.
Conversely a null polynomial has $D_d(R)$ annihilating every coefficient,
and the chosen regular element makes its rationalization zero.

Suppose $\mathcal N\ne0$. Select a nonzero coefficient $a$ of a null
polynomial of degree at most $d$. Then $D_d(R)a=0$. The monic polynomial
$P_d$ of B.064 gives a nonzero null polynomial $aP_d$. The tensor
\[
             aP_d\otimes P_d-P_d\otimes aP_d               \tag{B68.4}
\]
has zero polynomial image. It is nonzero: the two terms lie in distinct
summands $\mathcal N\otimes_R R[X]$ and $R[X]\otimes_R\mathcal N$ of the
tensor square of (B68.3), and extracting the top coefficient of the
ordinary polynomial $P_d$ sends the first term to $aP_d\ne0$.
Thus injectivity forces $\mathcal N=0$. The other implications now follow
from B.066, from (B68.3), and from the ordinary polynomial tensor
isomorphisms. This proves (B68.2).

For an infinite domain, each $D_d(R)$ is nonzero by choosing distinct
nodes. B.064 identifies (B68.1) with the stated ordinary
integer-valued-polynomial condition, proving that assertion.

For the polynomial examples, first assume every $D_d(A)$ contains a
nonzerodivisor of $A$, and put $R=A[Z]$. Such an element remains regular
in $R$, so it also belongs to $D_d(R)$ and proves the first hypothesis.
If $F\in\operatorname{Int}(R,Q(R))$ has degree $d$, evaluating on nodes
in $A$ and taking a finite $A$-linear combination of the corresponding
Cramer identities gives $v a_i\in A[Z]$ for some regular $v\in D_d(A)$.
Thus all coefficients belong to
\[
                      A[1/v][Z]\subset Q(R).
\]
Choose an integer $N$ strictly larger than the $Z$-degree of every
coefficient polynomial $a_i$. The degree intervals occupied by
$a_i(Z)Z^{Ni}$ are pairwise disjoint. Since $F(Z^N)\in A[Z]$, comparison
of these coefficients in the injective polynomial-ring inclusion
$A[Z]\subset A[1/v][Z]$ proves that every coefficient of every $a_i$
lies in $A$. Hence $F\in A[Z][X]$. B.064 gives (B68.1).

Finally, for arbitrary nonzero $A$, every $D_d(A[W])$ contains a monic
polynomial in $W$: for $d\ge2$ take the determinant on
$0,1,W,W^2,\ldots,W^{d-1}$, and for $d=1$ take $1$.
A monic polynomial over any nonzero commutative ring is a
nonzerodivisor. Applying the preceding argument with $A[W]$ in place
of $A$ proves the two-variable assertion. No localization of the
integer-valued module, flatness of $S$, or cancellation by zero divisors
has been used.
\end{proof}
